\subsection{Discussion}

The feature-representation ablation shows that self-supervised WavLM embeddings, used alone, provide a substantially more discriminative representation than handcrafted acoustic descriptors for this task: handcrafted features alone reached EER above 30\% for every model, while WavLM-only reduced this to below 27\% even at the smaller base checkpoint. Concatenating both feature families improved results only for Random Forest, whose per-tree feature subsampling appears more tolerant of a weaker auxiliary feature block than the boosting- and margin-based methods (XGBoost, SVM), for which concatenation measurably hurt performance relative to WavLM alone. This asymmetry across model families is itself a notable finding, and is treated as such rather than resolved into a single blanket recommendation.

The WavLM configuration and duration ablations together account for the largest improvements observed in this section. Moving from the base to the large WavLM checkpoint contributed the majority of the overall gain, with all-layer-averaged pooling and a fixed 3.0-second analysis window each contributing further, smaller but consistent improvements across all four classical models. The fixed 3.0-second window outperforming the duration range originally intended to match the anticipated 1--3 second inference scenario indicates that consistency and length of context matter more, for this feature-extraction approach, than duration diversity per se; this raises an open question, to be examined in the final report, of how a model trained exclusively on 3.0-second windows performs on genuinely short utterances at inference time, where shorter clips are extended by cyclic repetition rather than natural context.

The XGBoost regularization comparison confirms that an unregularized configuration substantially overfits the training partition, and that shallower trees with explicit L1/L2 penalties narrow the train--development gap considerably, though the comparison reported here is confounded with a concurrent feature and duration change and should be treated as directional rather than isolated evidence of the regularization effect alone.

With all eight baselines now trained on the full ASVspoof5 training partition and evaluated exactly once on the held-out evaluation partition (Section~6.4), the accuracy/efficiency trade-off motivating the inclusion of classical baselines can be assessed directly. The ranking by evaluation-set EER does not split cleanly along the deep-learning/classical-ML line motivating that comparison: NeXt-TDNN and SVM lead, followed by GBM, Random Forest, and XGBoost, with Spectro-AASIST, AASIST, and RawNet2 trailing. Notably, the two graph-attention and residual-recurrent architectures selected specifically for their reported state-of-the-art performance on ASVspoof benchmarks (AASIST and RawNet2) are, at present, the weakest baselines on ASVspoof5, well behind every classical model. Since this contradicts their reported literature performance by a wide margin, it is treated as an open implementation question rather than as evidence that the architectures themselves are unsuitable. For RawNet2, a newly captured development-set EER of 5.87\% (comparable to the classical baselines) makes under-training an unlikely explanation and points instead toward a generalization gap between the attack types seen during training and development and the 16 unseen attack algorithms introduced only in the evaluation partition (Section~5.1); whether the same explanation extends to AASIST, whose development-set EER has not yet been captured, remains open. Resolving both is a priority before the final report.

All results in this section were obtained on clean ASVspoof5 audio. No conclusions regarding RTC deployment readiness are drawn from these figures; RTC-induced degradation (the project's second objective) is measurable only once the RTC-aware augmentation pipeline is integrated into training and evaluation for the final report.